\appendix

\section*{Appendix:}

\section{Personas}

\subsection{Persona 1: Filip Philipsen}

\textbf{Quote:}  
\emph{“If I don’t fully understand it myself, I can’t explain it in a way that makes sense for the students.”}

\subsubsection{Personal Background}
Filip is a 50-year-old primary school teacher living in a small housing community in Bergen.
He has long experience working with children and parents and has spent most of his professional life in the educational sector.

\subsubsection{Professional Background}
Filip has 27 years of experience as a teacher and currently works at a primary school, where he teaches mathematics and physics.
His long tenure has given him deep insight into how students learn and how teaching conditions have changed over time.

\subsubsection{User Environment}
Filip teaches in a classroom setting with increasing use of digital tools, particularly after the COVID-19 pandemic.
He is required to continuously learn new teaching systems and works with students who experience more distractions
and reduced focus compared to earlier years.

\subsubsection{Goals}
\begin{itemize}
    \item Keep students motivated and engaged throughout lessons
    \item Help students experience mastery in the subjects
    \item Minimize failing and absenteeism in his class
\end{itemize}

\subsubsection{Psychographics}
Filip believes that learning methods must be adapted to individual students.
He prefers variation in teaching, such as short work periods combined with active breaks and practical tasks.
While he values repetition, he recognizes that homework does not work equally well for all students.
He is concerned about students’ increasing reliance on external motivation.

\subsubsection{Scenario}
When students struggle to stay focused or begin to fall behind, Filip wants to use varied,
practical teaching methods and clear structure to keep students engaged and help them genuinely understand the material.


% --------------------------------------------------

\subsection{Persona 2: Fiona Haaland}

\textbf{Quote:}  
\emph{“I want to understand what I’m studying and feel confident that my effort actually leads to good results.”}

\subsubsection{Personal Background}
Fiona is a 22-year-old university student living in a shared home with other students and her boyfriend in Bergen.
She recently transitioned from upper secondary school to university and finds the academic demands more independent
and challenging than before.

\subsubsection{Professional Background}
Fiona is currently a full-time student.

\subsubsection{Relevant Skills and Knowledge}
She has basic study skills such as note-taking and exam preparation, as well as basic programming knowledge
from upper secondary school.
She is familiar with standard digital tools commonly used in education.

\subsubsection{User Environment}
Fiona uses a laptop daily for schoolwork, assignments, studying, and gaming.
Her smartphone is mainly used for social media and entertainment, such as YouTube.
She primarily moves between home and university in her daily routine.

\subsubsection{Goals}
\textbf{Primary goal:}
\begin{itemize}
    \item Achieve good grades during the current semester
\end{itemize}

\textbf{Secondary and long-term goals:}
\begin{itemize}
    \item Complete her bachelor’s degree
    \item Obtain a stable job and provide for a future family
\end{itemize}

\subsubsection{Psychographics}
Fiona is motivated by a genuine desire to learn, but also by a fear of failure.
She values honesty and enjoys learning, often competing with herself to improve academically.

\subsubsection{Frustrations and Pain Points}
She becomes demotivated when she does not understand the material or when teachers appear unengaged or unpassionate about their subject.

\subsubsection{Personality Traits}
Fiona is introverted and observant, preferring to work independently rather than speak up in large classes.
She sets high expectations for herself and can become anxious when facing deadlines or difficult topics.

\subsubsection{Scenario}
When Fiona encounters a topic she does not understand during a lecture, she wants access to clear explanations and supportive learning resources so she can study independently, regain confidence, and stay motivated without feeling left behind.


% --------------------------------------------------

\subsection{Persona 3: Elias Eriksen}

\textbf{Quote:}  
\emph{“The school is scared to see me.”}

\subsubsection{Personal Background}
Elias is a 13-year-old student attending middle school (8th grade) in the Norwegian school system. He lives at home with his parents and has recently transitioned to middle school, where grades have started to matter more.

\subsubsection{Professional Background}
Elias is currently a middle school student.

\subsubsection{Relevant Personal Context}
Elias has struggled with his grades in the past and feels unmotivated toward school. The increased academic demands make him nervous, and he feels pressure to perform as grades become more important.

\subsubsection{User Environment}
Elias uses a smartphone and a laptop or computer. He mainly relies on digital platforms such as Itslearning and Wikipedia. His physical environment consists primarily of home and school.

\subsubsection{Goals}
\textbf{Primary goal:}
\begin{itemize}
    \item Achieve good grades
\end{itemize}

\textbf{Secondary and long-term goals:}
\begin{itemize}
    \item Get into the upper secondary school of his choice
    \item Discover what he wants to study in the future
\end{itemize}

\subsubsection{Psychographics}
Elias is motivated by a desire to avoid falling behind academically and seeks approval from teachers and parents. He wants to feel capable and successful at school.

\subsubsection{Attitudes and Values}
He believes school is important but finds it difficult to stay motivated. He values fairness and clear expectations and prefers practical explanations over abstract theory.

\subsubsection{Frustrations and Pain Points}
Elias struggles with increasingly complex subjects, low motivation, and difficulty focusing in class. He feels pressure as grades become more significant and is easily discouraged by poor results.

\subsubsection{Personality Traits}
Elias is somewhat insecure about his abilities, passive in classroom settings, and easily overwhelmed. He responds well to encouragement and clear structure.

\subsubsection{Scenario}
When school assignments become more difficult, Elias wants clear and simple help so that he can understand the material, keep up with his classmates, and improve his grades without feeling stressed or left behind.
